%% kapa: exam=1 round=21 year=2010 period=1 form=A subject=civil_law subject_ko=민법 seq=7 number=7
%% source: Q-Net 감정평가사 기출문제 (cst00309) artlSeq=5024181 file=extracted/1008/5024181/2010년 제21회 감정평가사 1차시험 기출문제/1교시.hwp
%% conversion: kordoc 4.24.51 -> tools/gen1.py (auto); review=proofread-sonnet; flags=none
\begin{question}{7}
甲은 남편 乙의 부동산을 몰래 담보로 제공하고 丙으로부터 금전을 차용하기로 마음먹었다. 甲과 공모한 丁은 위조된 乙의 신분증을 사용하여 자신이 乙인 것처럼 행세하면서, 丙으로부터 금전을 대출받고 담보목적으로 乙 소유의 부동산에 저당권을 설정해 주었다. 다음 설명으로 옳지 않은 것은? (다툼이 있으면 판례에 의함)
\begin{choices}
  \item 丁은 乙의 대리인이 아니다.
  \item 丙은 丁에게 대여금의 반환을 청구할 수 있다.
  \item 丙은 乙에게 권한을 넘은 표현대리책임을 물을 수 없다.
  \item 乙은 丙에게 저당권설정등기의 말소를 청구할 수 없다.
  \item 甲은 乙의 배우자로서 여전히 일상가사대리권을 보유한다.
\end{choices}
\end{question}
\answer{4}
